\section{Билет 31. Российская Федерация в начале XXI в.: перемены во внутренней политике и социально"=экономической сфере}

\subsection{Президент В.~В.~Путин и его окружение. Приоритеты нового руководства страны}

31 января 1999 г. "--- Президент РФ Б.~Н.~Ельцин во время новогоднего обращения к гражданам России объявил о своей отставке. Обязанности президента, в соответствии с Конституцией, возлагались на Председателя Правительства В.~В.~Путина, а новые президентские выборы назначались на май 2000 г.

Накануне в прессе появилась программная статья В.~Путина <<Россия на рубеже тысячелетий>>, где он излагал своё видение прошлого и представления о стоящих перед страной задачах. Основные положения статьи:

\begin{enumerate}
\item Ключ к возрождению и подъёму России находится в государственно"=политической сфере. Сильная государственная власть, в которой нуждается Россия "--- демократическое, правовое, дееспособное федеративное государство.
\item Необходимо повысить эффективность экономики, провести последовательную социальную политику, нацеленную на борьбу с бедностью и обеспечение устойчивого роста благосостояния населения.
\item Важна государственная поддержка науки, образования, культуры, здравоохранения.
\item Россия переживает один из самых трудных периодов своей истории и стоит перед лицом опасности <<оказаться во втором, а то и в третьем эшелоне государств мира>>.
\end{enumerate}

Май 2000 г. "--- В.~Путин был избран в первом туре выборов, набрав 52\% голосов избирателей.

\subsection{Окружение президента}

\subsubsection{Люди}
\begin{itemize}
\item Алексей Миллер "--- газовая сфера, государство при его руководстве получило контрольный пакет акций <<Газпрома>>;
\item Игорь Сечин "--- нефтяная отрасль, при нём государство вернуло активы <<Роснефти>>;
\item Владимир Якунин "--- железные дороги;
\item Сергей Чемезов "--- машиностроение, связь, оборонная промышленность;
\item Герман Греф "--- банковская сфера;
\item Сергей Лавров "--- министр иностранных дел;
\item Юрий Ковальчук "--- банки, страхование, медиа.
\end{itemize}

\subsubsection{Партии}
В Думе сформировалось устойчивое пропрезидентское большинство, стоящее на центристских позициях. В 1999 г. его составляли фракции <<Единство>>, <<Отечество "--- вся Россия>>, <<Народный депутат>> и группа <<Регионы России>>, которые вместе располагали 235 мандатами. Это позволяло проводить более активную политику, опираясь на поддержку законодателей. 

\subsection{Работа по совершенствованию общественных структур}
В 1990"~е годы в стране были заложены основы рыночной экономики и демократического устройства, для закрепления которых необходимо было навести порядок.

\subsubsection{Укрепление <<вертикали власти>>, создание федеральных округов}
Ослабление Центра не позволяло проводить в регионах России эффективный социально"=политический и экономический курс, поэтому насущной стала задача укрепления и усовершенствования <<вертикали власти>>.

Май 2000 г. "--- Указом Президента РФ на территории России создано 7 федеральных округов:
\begin{itemize}
\item Северо"=Западный;
\item Центральный;
\item Приволжский;
\item Южный;
\item Уральский;
\item Сибирский;
\item Дальневосточный.
\end{itemize}

Они не были ни новыми административными единицами, ни субъектами федерации. Их возглавили полномочные представители Президента, на которых возлагалась обязанность координировать деятельность местных органов власти на основе общероссийского законодательства, добиться неукоснительного исполнения Конституции всеми субъектами федерации. Развернулась большая работа по приведению региональных правовых актов в соответствие с общефедеральными нормами.

\subsubsection{Реформа Совета Федерации}
Была проведена реформа верхней палаты Федерального Собрания России "--- Совета Федерации, в который входили главы исполнительной власти (президенты, губернаторы) и законодательных собраний субъектов федерации. В 2000 г. места руководителей республик, краёв и областей стали постепенно занимать их представители, которые теперь работали в Совете Федерации на постоянной основе.

\subsubsection{Реформа российской армии}
Был запланирован рост финансирования ВС РФ при сокращении личного состава, оснащение войск современной боевой техникой. К 2010 г. предполагалось перейти преимущественно на контрактную основу комплектования армии. В июне 2002 г. был принят закон <<Об альтернативной гражданской службе>>.

\subsubsection{Перестройка системы МВД}
На базе МВД намечалось создать:

\begin{itemize}
\item муниципальную милицию (бытовые преступления);
\item федеральную полицию (угрозыск, ГАИ, борьба с наркотиками, лицензионно"=разрешительные и миграционные службы);
\item федеральную службу расследований;
\item Национальную гвардию.
\end{itemize}

\subsubsection{Другие изменения}
Произошли значительные структурные изменения в системе органов государственной безопасности. Произошла модернизация судебно"=правовой системы: в 2002 г. приняты новые акты, регулирующие деятельность суда и адвокатуры. В 2003 г. были подготовлены акты, определяющие основные параметры административной реформы.

\subsection{Экономическая политика в 2000--2004 гг.}
Министерство экономического развития и торговли подготовило обширный документ "--- <<Основные направления социально"=экономического развития РФ на долгосрочную перспективу>>, который определил стратегические приоритеты экономической политики государства на период до 2010 года.

\subsubsection{<<Равноудаление олигархов>> от власти}
Президент провозгласил необходимость ликвидации необоснованных привилегий и доходов, которые имели некоторые предприниматели за счёт особых отношений с чиновничеством и политиками. К этой теме общественное внимание было привлечено при расследовании <<дела ЮКОСА>> в 2003 году. Тогда были усилены акценты, касавшиеся ответственности крупного бизнеса перед обществом, и провозглашалось намерение создать равные <<правила игры>> для всех участников хоз. деятельности. Был принят закон РФ <<О противодействии отмыванию преступных доходов>>. В ноябре 2001 г. при Минфине был учреждён Комитет по финансовому мониторингу; Россия вступила в сообщество международных финансовых разведок.

\subsubsection{Реформы в налоговой сфере}
Конец 2000 г. "--- принятие новых Налогового и Таможенного кодекса.

Шкала налогов снижена с 30\% до 13\% и стала самой низкой в Европе, что способствовало <<выходу из тени>> многих секторов экономики. Были уменьшены таможенные ввозные пошлины.

\subsubsection{<<Естественные монополии>>}
Преобразования в управлении нефтяной, газовой индустрией, электроэнергетикой, в системе железнодорожного транспорта и сфере жилищно"=коммунального хозяйства направлены на создание конкурентной среды и увеличение инвестиций. Государственное управление полагалось неэффективным по определению, поэтому в собственности государства было решено оставить лишь объекты, необходимые для реализации его функций. На этой основе проводилась дальнейшая приватизация гос. предприятий.

\subsubsection{Земельная реформа}
Сентябрь 2001 г. "--- новый Земельный кодекс, закрепивший право собственности на землю и определивший механизм её купли"=продажи.

Июнь 2002 г. "--- закон <<Об обороте земель сельхозназначения>>, который санкционировал куплю"=продажу и этой категории земель.

На улучшение ситуации в аграрной сфере направлено льготное кредитование аграрных предприятий, увеличение поставок техники, некоторые протекционистские меры.

\subsubsection{Промышленность}
Государство не раз высказывало озабоченность насчёт <<сырьевого крена>> в развитии российской промышленности. В качестве стратегической поставлена задача добиться прогресса в отраслях, базирующихся на современных технологиях и производящих наукоёмкую продукцию. Государство стало уделять больше внимания ОПК (оборонно"=промышленному комплексу), его финансирование было расширено.

\subsubsection{Социальная направленность}
В течение 4 лет бюджет сводился с профицитом. Позитивные перемены в народном хозяйстве позволили значительно усилить социальную направленность проводимой политики.

К концу 2000"~го года удалось погасить большую часть задолженности по зарплате, пособиям и пенсиям. В 2001--2004 гг. несколько раз повышалась зарплата работникам бюджетной сферы, были увеличены пенсии. Появились новые рабочие места, что привело к сокращению числа безработных.

\subsection{Экономическая политика в 2004--2008 гг.}
К началу 2004 г. удалось добиться существенного снижения темпов инфляции: если в 1999 г. она составляла 36\%, то в 2003 г. "--- только 12\%. Рост ВВП составил в 2003 г. 7\%. Россия по объёму экономики вышла на седьмое место в мире.

С начала века товарооборот с зарубежными странами увеличился более чем в 5 раз. 2007 год был отмечен рекордным притоком капитала: он составил 82,3 млрд. долларов.

Август 2006 г. "--- досрочно был возвращён долг Парижскому клубу кредиторов.

По размеру золотовалютных резервов Россия заняла третье место в мире.

\subsection{Контртеррористическая операция на Северном Кавказе с 1999 г.}
23 сентября 1999 г. "--- ракетно"=бомбовому удару подвергнуты склады боеприпасов и оружия боевиков в окрестностях Грозного.

29 сентября 1999 г. "--- федеральные войска перешли границу с Чечнёй и начали наземные операции против террористов.

Февраль 2000 г. "--- федеральные войска взяли Грозный. Войсковая фаза операции была завершена, но в дальнейшем проявились сложности партизанской войны.

Зима 2001 г. "--- руководство операцией в Чечне передано от Министерства обороны Федеральной службе безопасности (Н.~П.~Патрушев).

Муфтий Чечни А.~Кадыров был назначен главой администрации. Утверждение в этом статусе бывшего активного участника первой чеченской войны свидетельствовало о готовности администрации РФ считаться со сложившимися реалиями.

Январь 2001 г. "--- в правительстве России была учреждена должность федерального министра по делам Чечни, создано правительство Чеченской республики, принята Федеральная программа восстановления её экономики и социальной сферы, на которую выделено 14,7 млрд. рублей.

23 октября 2002 г. "--- отряд чеченских террористов"=смертников под командованием М.~Бараева захватил в заложники около тысячи зрителей и актёров мюзикла <<Норд"=Ост>> в Театральном центре на Дубровке в Москве. Лидеры боевиков требовали от российского руководства полного вывода войск с территории Чечни, но в результате операции по спасению заложников почти всех террористов удалось уничтожить и предотвратить взрыв.

Март 2003 г. "--- был проведён референдум, в котором за принятие проекта Конституции высказалось 96\% голосовавших.

Октябрь 2003 г. "--- состоялись президентские и парламентские выборы. Началась подготовка договора о разграничении полномочий между федеральными и республиканскими органами власти, значительно увеличились ассигнования на восстановление разрушенной экономики и социальной сферы.

9 мая 2004 г. "--- жертвой теракта стал и президент Чеченской республики Ахмат Кадыров. 

\subsection{Президентство Д.~А.~Медведева}
14 марта 2004 г. В.~Путин на президентских выборах уже получил 71,2\% голосов и вступил в свой второй президентский срок. Партия <<Единая Россия>>, которая открыто заявляла о себе как о партии поддержки президента, одержала победу на выборах в Государственную Думу.

7 мая 2008 -- 7 мая 2012 гг. "--- президентство Дмитрия Медведева. Среди основных реформ, проведённых Медведевым, можно выделить следующие:

\begin{description}

\item[Модернизация] Президент назвал модернизацию вопросом выживания страны. По его мнению, необходимо было модернизировать экономику, производственную сферу, армию, медицину, технологии, образование и воспитание.

\item[Изменения конституции] Были внесены поправки об увеличении срока президентских полномочий с 4 до 6 лет, что стало первым изменением Конституции за её 15"=летнюю историю. Были также возвращены прямые выборы глав регионов, упрощён порядок регистрации политических партий.

\item[Удар по коррупции] Был подписан указ о создании Совета при Президенте РФ по противодействию коррупции. Медведев обязал госслужащих и членов их семей, а также руководителей госкорпораций и фондов предоставлять сведения о доходах и имуществе. В марте 2011 г. он потребовал вывода госчиновников из состава советов директоров крупных госкомпаний и госбанков.

\item[Договор о СНВ] В апреле 2010 г. главы РФ и США в Праге подписали новый договор о стратегических наступательных вооружениях, призванный стать одной из основ современной системы международной безопасности.

\item[Кавказский вопрос] В январе 2010 г. Медведев принял решение создать 8"=ой федеральный округ, в который вошли все северокавказские регионы. Центром нового округа стал Пятигорск.

\item[Часовые пояса] Было сокращено количество часовых поясов в России и изменено поясное время в ряде регионов.

\end{description}

\subsection{Итоги социально"=экономического, внутриполитического и внешнеполитического развития России к 2012 г.}

\subsubsection{Экономика}

\begin{enumerate}
\item 2001--2008 гг. "--- подъём экономики и её восстановление после кризиса. Прирост ВВП; переход от дефицита, неплатёжеспособности, зависимости от зарубежных кредитов и огромного гос. долга к профициту, значительным объёмам резервных фондов и одному из самых низких в мире уровней гос. долга в кратчайшие сроки.
\item 2009--2011 гг. "--- период мирового кризиса и посткризисного восстановления экономики. Успешная реализация Правительственной антикризисной программы, которая способствовала прекращению спада.
\item 2012 г. "--- переход к новой фазе роста, замедление как инвестиционного, так и потребительского спроса на фоне ослабления внешнего спроса.
\end{enumerate}

\subsubsection{Внутренняя политика}

\begin{itemize}
\item создание вертикали власти;
\item административная реформа "--- разграничение полномочий федеральных и местных органов власти, уменьшение бюрократического аппарата;
\item федеративная реформа "--- создание федеральных округов;
\item реформа парламентской системы;
\item победа на выборах провластной партии <<Единая Россия>>.
\end{itemize}

\subsubsection{Внешняя политика}

\begin{itemize}
\item восстановление Россией своих позиций на мировой сцене;
\item активное участие в международных организациях: БРИКС, ШОС, ООН, СНГ;
\item укрепление обороноспособности РФ;
\item стремление выстраивать многополярный мир;
\item сдерживание расширения НАТО на восток, попытка воспрепятствовать США вмешиваться в дела других государств.
\end{itemize}
